\section{Dynamic Sid-Only vs Dynamic Gather-Fix}

\subsection{区别是什么}

这两条线都属于 dynamic hint family，默认都用：

\begin{itemize}
\item \texttt{reward\_mode=rule\_only}
\item \texttt{dynamic\_hint\_max\_depth=3}
\item 不在 eval 侧继续应用 dynamic hint
\end{itemize}

真正的差异主要有两点：

\begin{itemize}
\item \texttt{dynamic sid-only} 不是在 mixed RL 数据上只改一个 flag，而是直接切到了专门的 sid-only RL 数据变体：\path{Instruments_grec_rlsidonly_index_emb-...}。
\item \texttt{dynamic gather-fix} 则继续跑默认 mixed RL 数据变体：\path{Instruments_grec_index_emb-...}；在本文里，它也被视为当前 canonical dynamic 实现。
\end{itemize}

\subsection{对应脚本}

{\footnotesize
\begin{itemize}
\item 画图脚本：\path{docs/research/scripts/build_instruments_report_figures.py}
\item 脚本目录：\path{hope/Qwen2_5-3B-Isntruct-qwen4B-4-256-MIMIGenRec-grec/}
\item canonical dynamic：\path{Qwen2_5-3B-Isntruct-qwen4B-4-256-MIMIGenRec-grec-rl-rule-only-dynamic-hint.sh}
\item dynamic sid-only：\path{Qwen2_5-3B-Isntruct-qwen4B-4-256-MIMIGenRec-grec-rl-rule-only-dynamic-hint-sid-only.sh}
\end{itemize}
}

\subsection{关键指标}

\begin{table}[H]
\centering
\small
\setlength{\tabcolsep}{4pt}
\begin{tabular}{p{4.0cm} p{2.5cm} c c c c c}
\toprule
Variant & Best ckpt & Epoch & NDCG@10 & HR@10 & NDCG@50 & HR@50 \\
\midrule
dynamic sid-only & checkpoint-2394 & 1.805 & 0.0921 & 0.1155 & 0.1065 & 0.1830 \\
dynamic gather-fix & checkpoint-2997 & 1.802 & 0.0936 & 0.1169 & 0.1083 & 0.1855 \\
\bottomrule
\end{tabular}
\caption{dynamic sid-only 与 canonical dynamic gather-fix。}
\end{table}

\subsection{图表}

\begin{figure}[H]
\centering
\includegraphics[width=0.96\textwidth]{dynamic_sid_only_vs_dynamic_gather_fix_curves.png}
\caption{dynamic sid-only vs canonical dynamic gather-fix。}
\end{figure}

\subsection{结论}

\begin{itemize}
\item \texttt{dynamic sid-only} 不是完全失败，但 peak 来得更早，后续没有继续抬高。
\item \texttt{dynamic gather-fix} 四个主指标同时更强，因此应被视为当前默认的 dynamic 实现。
\end{itemize}
